\documentclass[11pt]{article}

\usepackage[margin=2.5cm]{geometry}
\usepackage{amsmath}
\usepackage{amssymb}
\usepackage{booktabs}
\usepackage{array}

\title{Sistema de probabilidad usado en \texttt{zanni\_pytorch.py}}
\author{}
\date{}

\begin{document}

\maketitle

\section{Resumen}

El script \texttt{zanni\_pytorch.py} calcula una estrategia optima de ofertas
para el problema ``Zanni de Oro''. Aunque el nombre del archivo menciona
\texttt{pytorch}, la version actual no usa PyTorch ni librerias externas: usa un
modelo probabilistico continuo, condiciones de primer orden y busqueda binaria.

El intervalo de precios usado por el programa es
\[
  [L,H] = [620,880],
\]
por lo que su largo es
\[
  S = H-L = 260.
\]
El codigo tambien considera
\[
  N = 10000
\]
monedas.

\section{Normalizacion}

Para trabajar con formulas mas simples, cada oferta real \(b\) se transforma a
una oferta normalizada \(t\) en el intervalo \([0,1]\):
\[
  t = \frac{b-L}{S}.
\]
Equivalentemente, una oferta normalizada \(t\) corresponde al precio real
\[
  b = L + S t.
\]

Asi, \(t=0\) representa \(620\), y \(t=1\) representa \(880\).

\section{Distribucion probabilistica}

El programa modela el precio de reserva normalizado \(R\) con funcion de
distribucion acumulada
\[
  F(t) = t^2, \qquad 0 \le t \le 1.
\]
Por lo tanto,
\[
  \Pr(a < R \le b) = F(b)-F(a) = b^2-a^2.
\]

La densidad asociada es
\[
  f(t) = F'(t) = 2t.
\]
Esto significa que los valores altos de reserva son mas probables que los bajos,
porque la densidad crece linealmente con \(t\).

\section{Ganancia esperada}

Supongamos que se hacen \(k\) ofertas normalizadas
\[
  0 < t_1 < t_2 < \cdots < t_k \le 1.
\]
Definimos \(t_0=0\). Si la reserva cae en el intervalo
\[
  t_{i-1} < R \le t_i,
\]
entonces la oferta que se acepta es \(t_i\). La ganancia normalizada de esa
oferta es
\[
  1-t_i,
\]
y la probabilidad de que ocurra ese caso es
\[
  \Pr(t_{i-1}<R\le t_i)=t_i^2-t_{i-1}^2.
\]

Por lo tanto, la ganancia esperada normalizada es
\[
  G(t_1,\ldots,t_k)
  =
  \sum_{i=1}^{k}
  (1-t_i)(t_i^2-t_{i-1}^2).
\]

El programa implementa exactamente esta expresion en la funcion
\[
  \texttt{expected\_profit\_normalized}.
\]

\section{Optimizacion}

Para maximizar \(G\), el codigo usa las condiciones estacionarias. Para una
oferta interna \(t_i\), con \(1 \le i < k\), la condicion de primer orden lleva a
\[
  t_{i-1}^2 - 3t_i^2 + 2t_i t_{i+1}=0.
\]
Despejando \(t_{i+1}\), se obtiene la recurrencia usada por el programa:
\[
  t_{i+1}
  =
  \frac{3t_i^2-t_{i-1}^2}{2t_i}.
\]

La ultima oferta \(t_k\) satisface la condicion terminal
\[
  t_{k-1}^2 + 2t_k - 3t_k^2 = 0.
\]
Para el caso \(k=1\), esta condicion queda como
\[
  2t_1 - 3t_1^2 = 0,
\]
de donde el optimo no trivial es
\[
  t_1=\frac{2}{3}.
\]

El algoritmo elige \(t_1\) mediante busqueda binaria en \([0,1]\), genera las
siguientes ofertas con la recurrencia anterior y ajusta \(t_1\) hasta cumplir la
condicion terminal.

\section{Conversion a ganancias reales}

Una vez obtenida la ganancia esperada normalizada \(G\), el programa calcula la
ganancia esperada por moneda como
\[
  \text{ganancia por moneda} = S \cdot G.
\]

Como hay \(N=10000\) monedas, la ganancia esperada total es
\[
  \text{ganancia total} = N \cdot S \cdot G.
\]

\section{Resultados del script}

Ejecutando
\[
  \texttt{python3 zanni\_pytorch.py --max-bids 5}
\]
se obtienen los siguientes valores:

\begin{center}
\begin{tabular}{c>{\raggedright\arraybackslash}p{7.3cm}rr}
\toprule
Fase & Ofertas optimas & Ganancia por moneda & Ganancia total \\
\midrule
1 & 793.333333 & 38.518519 & 385185.19 \\
2 & 755.652174,\; 823.478261 & 53.081285 & 530812.85 \\
3 & 733.634204,\; 790.451306,\; 837.798892 & 60.815843 & 608158.43 \\
4 & 718.886355,\; 768.329533,\; 809.532181,\; 846.256281 & 65.630647 & 656306.47 \\
5 & 708.184178,\; 752.276267,\; 789.019674,\; 821.769233,\; 851.860970 & 68.922448 & 689224.48 \\
\bottomrule
\end{tabular}
\end{center}

\section{Pseudocodigo}

\[
\begin{array}{ll}
\textbf{entrada:} & k \text{ numero de ofertas} \\
\textbf{salida:} & t_1,\ldots,t_k \text{ ofertas normalizadas optimas}
\end{array}
\]

\begin{enumerate}
  \item Buscar \(t_1\) en el intervalo \([0,1]\).
  \item Fijar \(t_0=0\).
  \item Para \(i=1,\ldots,k-1\), calcular
  \[
    t_{i+1}=\frac{3t_i^2-t_{i-1}^2}{2t_i}.
  \]
  \item Evaluar la condicion terminal
  \[
    t_{k-1}^2 + 2t_k - 3t_k^2.
  \]
  \item Ajustar \(t_1\) por busqueda binaria hasta que la condicion terminal sea
  cercana a cero.
  \item Convertir cada \(t_i\) a precio real con
  \[
    b_i=620+260t_i.
  \]
\end{enumerate}

\section{Conclusion}

El sistema probabilistico usado por el programa se basa en una reserva aleatoria
continua con distribucion \(F(t)=t^2\). La estrategia optima no se obtiene por
simulacion, sino maximizando la ganancia esperada mediante condiciones
analiticas. La unica parte numerica es la busqueda binaria del primer punto
\(t_1\), que permite construir todas las demas ofertas.

\end{document}
